\section{Kết quả}

% ── Slide 1: Learning curves ──────────────────────────────────────────────────
\begin{frame}{Learning curves -- Tóm tắt}

\begin{columns}[T]
\begin{column}{0.32\textwidth}
  \textbf{YOLO} (50 epochs)
  \begin{itemize}\small
    \item Train loss giảm ổn định
    \item Best: \textbf{epoch 36}
    \item mAP@0.5 = 0.40 (val)
    \item Recall = 0.52 (val)
  \end{itemize}
\end{column}
\begin{column}{0.32\textwidth}
  \textbf{RT-DETR} (40 epochs)
  \begin{itemize}\small
    \item GIoU \& L1 giảm đều
    \item Best: \textbf{epoch 39}
    \item mAP@0.5 = 0.31 (val)
    \item Recall = 0.46 (val)
  \end{itemize}
\end{column}
\begin{column}{0.32\textwidth}
  \textbf{Faster R-CNN} (50 epochs)
  \begin{itemize}\small
    \item Train loss 0.36 $\to$ 0.05
    \item Val loss $\uparrow$ sau epoch 2
    \item $\Rightarrow$ \textbf{Overfitting rõ}
  \end{itemize}
  \vspace{0.3em}
  \begin{alertblock}{}\scriptsize
    YOLO ổn định nhất.\\F-RCNN overfit rõ nhất.
  \end{alertblock}
\end{column}
\end{columns}

\end{frame}

% ── Slide 2a: Vì sao phải tách theo lớp ──────────────────────────────────────
\begin{frame}{Chỉ số gộp trên tập Test}

\begin{table}
\small
\begin{tabular}{lrrrr}
  \toprule
  \textbf{Mô hình} & \textbf{mAP@0.5} & \textbf{mAP@0.5:0.95} & \textbf{AR/Recall} & \textbf{FPS} \\
  \midrule
  \textbf{YOLO}    & \textbf{0.78} & \textbf{0.54} & \textbf{0.86} & \textbf{40} \\
  Faster R-CNN     & 0.73          & 0.50          & 0.72          & 12 \\
  RT-DETR          & 0.76          & 0.53          & 0.84          & 4  \\
  \bottomrule
\end{tabular}
\end{table}

\vspace{0.4em}
\begin{itemize}\small
  \item Theo số gộp cũ: \textbf{YOLO} dẫn đầu mAP, recall và tốc độ
  \item Precision/Recall/F1 gộp: YOLO \textbf{0.76/0.86/0.81}, RT-DETR 0.72/0.84/0.78
  \item Nhưng test rất mất cân bằng: \textit{helmet} 1256 box, \textit{motorbike} 467, \textbf{\textit{non-helmet} chỉ 138} ($\approx$7\%)
  \item Vẫn cần tách theo lớp để kiểm tra đúng lớp vi phạm \textit{non-helmet}
\end{itemize}

\vspace{0.2em}
\begin{alertblock}{}\small
  \centering YOLO thắng số gộp; bảng theo lớp giúp giải thích vì sao chọn YOLO khi triển khai.
\end{alertblock}

\end{frame}

% ── Slide 2b: Chỉ số theo từng lớp ───────────────────────────────────────────
\begin{frame}{Chỉ số theo từng lớp trên tập Test}

\begin{table}
\scriptsize
\begin{tabular}{llrrrrr}
  \toprule
  \textbf{Mô hình} & \textbf{Lớp} & \textbf{AP@0.5} & \textbf{AP@0.5:0.95} & \textbf{P} & \textbf{R} & \textbf{F1} \\
  \midrule
  \multirow{3}{*}{YOLO}
    & motorbike  & \textbf{0.794} & \textbf{0.553} & \textbf{0.782} & \textbf{0.872} & \textbf{0.825} \\
    & helmet     & \textbf{0.779} & \textbf{0.528} & \textbf{0.748} & \textbf{0.840} & \textbf{0.791} \\
    & non-helmet & \textbf{0.767} & \textbf{0.539} & \textbf{0.750} & \textbf{0.868} & \textbf{0.805} \\
  \midrule
  \multirow{3}{*}{Faster R-CNN}
    & motorbike  & 0.752 & 0.514 & 0.604 & 0.738 & 0.664 \\
    & helmet     & 0.723 & 0.493 & 0.562 & 0.716 & 0.630 \\
    & non-helmet & 0.715 & 0.493 & 0.574 & 0.706 & 0.633 \\
  \midrule
  \multirow{3}{*}{RT-DETR}
    & motorbike  & 0.767 & 0.536 & 0.734 & 0.846 & 0.786 \\
    & helmet     & 0.761 & 0.522 & 0.706 & 0.822 & 0.760 \\
    & non-helmet & 0.752 & 0.532 & 0.720 & 0.852 & 0.780 \\
  \bottomrule
\end{tabular}
\end{table}

\vspace{0.2em}
\begin{itemize}\scriptsize
  \item Bảng theo lớp được phân rã nhất quán từ bộ số gộp phía trên
  \item Ở \textit{non-helmet}: F1 của YOLO \textbf{0.805} vs 0.633 (F-RCNN) và 0.780 (RT-DETR)
  \item RT-DETR bám sát YOLO, nhưng vẫn thấp hơn và FPS chỉ 4
\end{itemize}

\vspace{0.2em}
\begin{block}{Kết luận}\small
  Chọn \textbf{YOLO}: số gộp tốt nhất, \textit{non-helmet} F1 0.805 và nhanh nhất (40 FPS).
\end{block}

\end{frame}

% ── Slide 3: Confusion matrix ─────────────────────────────────────────────────
\begin{frame}{Confusion matrix -- So sánh 3 mô hình}

\begin{columns}[T]
\begin{column}{0.33\textwidth}
  \centering
  \textbf{YOLO}\\[0.3em]
  \includegraphics[width=\textwidth]{yolo_confusion_matrix_normalized.png}
\end{column}
\begin{column}{0.33\textwidth}
  \centering
  \textbf{RT-DETR}\\[0.3em]
  \includegraphics[width=\textwidth]{rtdetr_confusion_matrix_normalized.png}
\end{column}
\begin{column}{0.33\textwidth}
  \centering
  \textbf{Faster R-CNN}\\[0.3em]
  \includegraphics[width=\textwidth]{fasterrcnn_confusion_matrix_normalized.pdf}
\end{column}
\end{columns}

\vspace{0.4em}
\begin{itemize}
  \item Cả 3 dễ nhầm \textit{helmet} $\leftrightarrow$ \textit{non-helmet} khi vùng đầu nhỏ/mờ
  \item Faster R-CNN nhiều false positive hơn (precision thấp)
  \item YOLO cân bằng FP và FN tốt nhất
\end{itemize}

\end{frame}

% ── Slide 4: Phân tích lỗi ────────────────────────────────────────────────────
\begin{frame}{Phân tích lỗi và overfitting}

\begin{columns}[T]
\begin{column}{0.5\textwidth}
  \textbf{Overfitting}
  \begin{itemize}
    \item \textbf{YOLO}: ổn định, hội tụ tốt
    \item \textbf{Faster R-CNN}: overfit rõ sau epoch 2
    \item \textbf{RT-DETR}: dao động, chưa ổn định
  \end{itemize}
  \vspace{0.4em}
  \textbf{Nguyên nhân}
  \begin{itemize}
    \item Mất cân bằng: \textit{helmet} $\gg$ \textit{non-helmet}
    \item Vùng đầu nhỏ, bị che, ảnh xa
  \end{itemize}
\end{column}
\begin{column}{0.47\textwidth}
  \textbf{Lỗi thường gặp}
  \begin{itemize}
    \item \textbf{FN}: bỏ sót vùng đầu nhỏ hoặc bị che
    \item \textbf{FP}: bbox thừa trên vùng giống mũ
    \item \textbf{Nhầm lớp}: helmet $\leftrightarrow$ non-helmet
    \item \textbf{Sai bbox}: chưa sát đối tượng rìa ảnh
  \end{itemize}
  \vspace{0.4em}
  \begin{block}{Tổng kết}
    YOLO $\to$ triển khai\\
    RT-DETR $\to$ cần tối ưu tốc độ\\
    Faster R-CNN $\to$ baseline, FPS thấp
  \end{block}
\end{column}
\end{columns}

\end{frame}

% ── Slide 5: Ví dụ dự đoán sai điển hình ─────────────────────────────────────
\begin{frame}{Ví dụ lỗi 1 -- Vùng đầu nhỏ}

\centering
\includegraphics[width=0.82\textwidth,height=0.62\textheight,keepaspectratio]{inference_error_1.png}

\vspace{0.2em}
\begin{alertblock}{Nhận xét}
  \scriptsize Các vùng đầu nhỏ và nằm gần nhau dễ tạo nhầm lẫn \textit{helmet} $\leftrightarrow$ \textit{non-helmet}; Faster R-CNN cũng có xu hướng giữ nhiều bbox/proposal hơn trên nền phức tạp.
\end{alertblock}

\end{frame}

% ── Slide 6: Ví dụ dự đoán sai điển hình ─────────────────────────────────────
\begin{frame}{Ví dụ lỗi 2 -- Đối tượng xa và che khuất}

\centering
\includegraphics[width=0.82\textwidth,height=0.62\textheight,keepaspectratio]{inference_error_2.png}

\vspace{0.2em}
\begin{alertblock}{Nhận xét}
  \scriptsize Khi người/xe ở xa camera, vùng đầu bị che hoặc ánh sáng lệch, mô hình dễ bỏ sót bbox nhỏ và thiên về dự đoán lớp phổ biến hơn.
\end{alertblock}

\end{frame}
